\section*{七、思考題}
\textbf{Q1. 第一輪分錯的樣本，為何第二輪容易重複抽到？}

ID 8、9 的權重由 0.1 升至 0.25，每筆是分對樣本權重 0.0625 的四倍；若依此機率有放回抽樣，就較容易重複抽到它們。
本例第二輪的固定清單中，ID 8 出現兩次、ID 9 出現三次，但權重較高不保證每次抽樣都會重複。

\textbf{Q2. 為何計算 $\epsilon_2$ 仍要回到原始資料及其權重？}

抽樣資料只用來選規則，$\epsilon_2$ 的定義則是原始 10 筆資料中錯分樣本的權重總和。
本例規則 E 的抽樣錯誤率為 $1/10$，但原始資料錯分 ID 5、6、7，故 $\epsilon_2=3/16$。

\textbf{Q3. 為何第二輪可能選出不同切分規則？}

抽樣組成改變後，各規則的未加權錯誤筆數也隨之改變；第一輪 A 錯 2 筆、E 錯 3 筆，所以選 A。
第二輪重複抽到 ID 8、9，使 A 錯 5 筆，而 E 只錯 1 筆，因此改選切分點 9.5 的 E。
